\documentclass[tikz,border=8pt]{standalone}
\usepackage{ctex}
\usepackage{amsmath}
\usetikzlibrary{calc, arrows.meta}

\begin{document}
% 面面位置关系：平行 / 相交
\begin{tikzpicture}[
  x={(1cm,0cm)},
  y={(0.38cm,0.28cm)},
  z={(0cm,0.85cm)},
  thick,
  plane1/.style={fill=blue!12, draw=blue!50, fill opacity=0.75},
  plane2/.style={fill=green!12, draw=green!60!black, fill opacity=0.75},
]

%%% 栏1：两平面平行 α ∥ β
\begin{scope}
  \node[font=\small\bfseries] at (1.5, 4.6, 0) {(1) 两平面平行};
  \node[font=\small] at (1.5, 4.2, 0) {$\alpha \parallel \beta$};

  % 下方平面 β（较低）
  \fill[plane2] (0,0,0) -- (3,0,0) -- (3.8,2,0) -- (0.8,2,0) -- cycle;
  \node[green!60!black, font=\small] at (4.1, 1.0, 0) {$\beta$};

  % 上方平面 α（较高，平行偏移）
  \fill[plane1] (0,0,1.6) -- (3,0,1.6) -- (3.8,2,1.6) -- (0.8,2,1.6) -- cycle;
  \node[blue!70, font=\small] at (4.1, 1.0, 1.6) {$\alpha$};

  % 平行虚线（两平面对应点连线）
  \draw[gray, dashed, thin] (0,0,0) -- (0,0,1.6);
  \draw[gray, dashed, thin] (3,0,0) -- (3,0,1.6);
  \draw[gray, dashed, thin] (0.8,2,0) -- (0.8,2,1.6);
\end{scope}

%%% 栏2：两平面相交 α ∩ β = l
\begin{scope}[xshift=6.5cm]
  \node[font=\small\bfseries] at (1.5, 4.6, 0) {(2) 两平面相交};
  \node[font=\small] at (1.5, 4.2, 0) {$\alpha \cap \beta = l$};

  % 交线 l（垂直于观察方向，偏中）
  % 平面 α（左倾）
  \fill[plane1]
    (0.4,2,0) -- (3.4,2,0) -- (1.8,2,-1.2) -- (-0.8,2,-1.2)
    -- (1.2,2,1.6) -- cycle;
  % 更简洁的平面 α（右上斜面）
  \begin{scope}
    % 交线：从(0,1,0)到(3,1,0)，即y=1平面
    % 平面 α：过交线，往上方偏右展开
    \fill[plane1, opacity=0.65]
      (0,1,0) -- (3,1,0) -- (3.5,2.5,1.6) -- (0.5,0.5,1.6) -- cycle;
    \node[blue!70, font=\small] at (3.8, 1.8, 0.9) {$\alpha$};

    % 平面 β：过交线，往下方偏左展开
    \fill[plane2, opacity=0.65]
      (0,1,0) -- (3,1,0) -- (3.8,2.8,-1.0) -- (0.8,0.2,-1.0) -- cycle;
    \node[green!60!black, font=\small] at (4.0, 1.8, -0.6) {$\beta$};

    % 交线（红色加粗）
    \draw[red, thick] (0,1,0) -- (3,1,0);
    \node[red, font=\small, above] at (1.5,1,0) {$l$};

    % 端点标记
    \fill[red] (0,1,0) circle [radius=1.5pt];
    \fill[red] (3,1,0) circle [radius=1.5pt];
  \end{scope}
\end{scope}

\end{tikzpicture}
\end{document}
